\uuid{n6kf}
\titre{Calcul d'intégrales définies}
\niveau{L2}
\module{Analyse}
\chapitre{Calcul d'intégrales}
\sousChapitre{Primitives}
\cpge{mpsi-mp2i.primitives-equations-differentielles.a, pcsi.techniques-calcul-primitives.a, ptsi.techniques-calcul-primitives.a, mpsi-mp2i.integration.c, pcsi.integration.b, ptsi.integration.b, bcpst1.derivees-primitives-integrales.d}
\theme{polynôme, puissance, fraction rationnelle}
\auteur{Erwan L'Haridon}
\datecreate{2026-02-10}
\organisation{AMSCC}
\difficulte{2}

\contenu{
    \texte{Calculer les intégrales suivantes :}
    \begin{enumerate}
        \item \question{\( I_1 = \int_{-1}^{1} \left(t^4 + t^{917}\right) dt \)}
        \item \question{\( I_2 = \int_{1}^{2} s^{\frac{4}{3}} ds \)}
        \item \question{\( I_3 = \int_{1}^{2} \frac{dt}{(t + 4)^3} \)}
        \item \question{\( I_4 = \int_{1}^{2} \frac{\left(\sqrt{x}\right)^3}{x^4} \cdot \frac{1}{\sqrt[3]{x^2}} dx \)}
    \end{enumerate}
}
